\section{引言与文献综述}

\subsection{研究背景}

在教育数字化转型背景下，高校地理科学专业承担着培养具备GIS与AI素养的复合型专业人才的使命。地理空间人工智能（GeoAI）正深刻改变地理数据的分析与应用方式\citep{Jiang2023GeoAI}，GIS空间分析能力已成为地理专业人才的核心竞争力。

然而，传统GIS实践教学存在明显痛点：教学内容偏重软件操作，学生"会点菜单，不会建模"\citep{杜清运2017}；空间分析过程抽象，缺乏真实项目驱动；教学内容与科研实践脱节\citep{王民2018}。生成式AI的快速发展为解决上述问题提供了新思路——大语言模型可作为"智能协作者"辅助代码生成、模型构建与结果解读\citep{Dwivedi2023}。

\subsection{文献综述}

近年来，AI赋能GIS教育已成为研究热点。\citet{Zhang2023AI}研究表明，AI辅助工具可显著提升GIS学习效率与深度；\citet{Wang2020}指出AI语言模型在GIS教育中具有应用潜力；\citet{Zhao2021}探索了人工智能赋能地理实践教学的模式与路径；\citet{Luo2022}提出了人机协同视域下的地理智慧教育理论框架。在国际领域，GeoAI教育研究聚焦于如何将机器学习与空间分析有机整合\citep{Klosterman2015}，以及GIS空间分析的教学内容与方法改革\citep{杜清运2017}。

已有研究存在以下不足：（1）缺乏针对AI与GIS实践深度融合的可操作性教学设计；（2）评价体系研究薄弱，尤其是AI辅助教学的过程性评价；（3）较少关注AI时代地理人才核心能力的重新定义。本研究针对上述不足，设计了基于生成式AI的城市物流园区选址分析教学案例，构建了四阶段人机协同教学模式。

\subsection{研究贡献}

本文贡献包括：（1）构建"问题定义—数据获取—空间建模—成果表达"四阶段人机协同教学模式；（2）设计针对高校地理专业特点的AI辅助GIS教学流程与Prompt示例；（3）建立突出过程性评价的多元化教学评价体系。全文组织如下：第2节介绍教学目标；第3节阐述案例设计；第4节论述四阶段教学过程；第5节构建评价体系；第6节总结创新点。